% =============================================================================
% NeurIPS 2026 — ANONYMIZED SUBMISSION
%
% REQUIRED: place the official neurips_2026.sty next to this file.
%   Download the author kit from
%   https://neurips.cc/Conferences/2026/CallForPapers
%   (or open the official Overleaf template and export the source).
%   Do NOT use any third-party reconstruction of the style file.
%
% The DEFAULT option is the anonymous double-blind submission mode.
% Do not pass [preprint] or [final] for review copies.
%
% BEFORE SUBMITTING:
%   1. Replace ANONYMIZED-REPO-URL with an anonymous.4open.science link
%   2. Confirm the workshop's page limit (main track is 9 content pages;
%      references, checklist, and appendices do not count)
%   3. Compile and read the PDF for any identifying string
% =============================================================================
\documentclass{article}

\PassOptionsToPackage{numbers,sort&compress}{natbib}

% Workshop submission, double-blind. The style file provides dedicated
% workshop options; the bare default is the MAIN TRACK, not a workshop.
%   [dblblindworkshop] -> anonymous (use this unless told otherwise)
%   [sglblindworkshop] -> NOT anonymous; only if the workshop is single-blind
\usepackage[dblblindworkshop]{neurips_2026}
\workshoptitle{SLM-Agents: 1st NeurIPS Workshop on SLMs for Agentic Systems}

\usepackage[utf8]{inputenc}
\usepackage[T1]{fontenc}
\usepackage{hyperref}
\usepackage{url}
\usepackage{booktabs}
\usepackage{amsmath}
\usepackage{graphicx}
\usepackage{multirow}
\usepackage{xcolor}        % required by the checklist answer macros
\usepackage{caption}
\captionsetup{font=small}

\title{Quantized Small Language Models as Tool-Using Agents:\\
Memory Savings for Free, Fine-Tuning with Care}

% Anonymous for review. The style file suppresses author details in the
% default (submission) mode; the block below is the template's standard form.
\author{%
  Anonymous Author(s)\\
  Affiliation\\
  Address\\
  \texttt{email}
}

\begin{document}
\maketitle

\begin{abstract}
Agentic AI systems typically rely on large language models whose memory
footprint prevents deployment on constrained hardware. We present a
systematic empirical study of whether a \emph{small} language model
(Phi-3-mini, 3.8B parameters) remains a capable tool-using agent under
NF4 4-bit quantization, and whether lightweight QLoRA fine-tuning on
generic agent demonstrations helps or hurts. Using a single unified ReAct
evaluation harness across three benchmarks---a 50-task diagnostic agentic
suite, GSM8K with a calculator tool, and HotpotQA (distractor) with a
search tool---we find that INT4 quantization is essentially free for
agentic capability. On the diagnostic suite the FP16 and INT4 agents
succeed and fail on \emph{exactly the same tasks} (zero discordant pairs,
$n{=}50$); on the standard benchmarks paired McNemar tests find no
significant difference ($p{=}0.58$ on GSM8K, $p{=}0.63$ on HotpotQA),
with weight memory falling from 7.64\,GB to 2.21\,GB (71\%). We document
two deployment-relevant caveats. First, on a T4, INT4 inference is
roughly $2\times$ \emph{slower} than FP16 because dequantization
overhead dominates at batch size one, so quantization buys memory, not
speed, on commodity accelerators. Second, QLoRA fine-tuning on 15 short
generic ReAct demonstrations (0.25\% of parameters, under two minutes of
training) leaves diagnostic-suite accuracy unchanged but degrades
performance on the standard benchmarks (pooled McNemar $p{=}0.041$),
concentrated on multi-hop question answering. Trace analysis identifies the
mechanism unambiguously: the fine-tuned agent issues a second search on 10
of 100 questions where the base INT4 agent issues one on 39 (Wilcoxon
signed-rank $p{<}10^{-5}$), and second searches are precisely the ones that
pay off. Our results
suggest quantized SLM agents are viable for deployment, but that naive
behavioral fine-tuning on short traces can silently truncate exploratory
behavior in a way that per-benchmark accuracy alone does not reveal.
Code, evaluation harness, per-task traces, and all raw results are
released.
\end{abstract}

\section{Introduction}
\label{sec:intro}

Agentic AI systems---language models that decompose tasks, invoke external
tools, and act over multiple steps---have advanced rapidly, but almost
exclusively on top of large proprietary models. This creates a deployment
gap: the settings where autonomous agents are most valuable (on-device
assistants, on-premise enterprise automation, low-cost cloud inference)
are precisely those where serving a 70B+ parameter model is infeasible.
Recent position work argues that small language models (SLMs) are the
natural substrate for agentic AI~\cite{belcak2025slm}, since agent
workloads are dominated by narrow, repetitive tool-calling rather than
open-ended generation.

Two efficiency levers make SLM agents plausible: post-training
quantization~\cite{dettmers2023qlora,dettmers2022llmint8}, which shrinks
memory at some risk to capability, and parameter-efficient fine-tuning
(PEFT)~\cite{hu2022lora,dettmers2023qlora}, which can specialize behavior
at negligible cost. What is missing is a careful account of how these
levers interact with \emph{agentic} behavior specifically: multi-step
reasoning, tool selection, and iterative search are plausibly more
fragile under quantization than single-turn generation, and behavioral
fine-tuning may have out-of-distribution side effects that single-metric
evaluations miss.

This paper contributes such an account for one representative open SLM,
Phi-3-mini (3.8B)~\cite{abdin2024phi3}, evaluated end-to-end as a ReAct
agent~\cite{yao2023react} on free commodity hardware (a single NVIDIA
T4). We ask three questions. \textbf{(Q1)} How much agentic capability
is lost at INT4 precision? \textbf{(Q2)} What are the true latency and
memory trade-offs on commodity GPUs? \textbf{(Q3)} Does lightweight QLoRA
fine-tuning on a handful of generic agent demonstrations improve agentic
behavior, and does any change transfer to standard benchmarks?

Because every configuration runs the same task list in the same order
under a single harness, our results are \emph{paired}. We therefore use
McNemar's exact test for accuracy and Wilcoxon signed-rank for tool-call
depth throughout, which is substantially more informative at $n{=}100$
than comparing independent proportions. Our findings:

\begin{enumerate}
  \item \textbf{INT4 quantization is essentially free for agentic
  capability.} On the diagnostic suite, FP16 and INT4 produce
  \emph{identical} per-task outcomes---zero discordant pairs across 50
  tasks. On GSM8K and HotpotQA, paired tests find no significant
  difference ($p{=}0.58$, $p{=}0.63$), and the paired confidence
  intervals bound any INT4 penalty at roughly 11 accuracy points.
  Weight memory falls 71\% (7.64\,GB\,$\rightarrow$\,2.21\,GB) (Q1).
  \item \textbf{Quantization buys memory, not speed, on commodity
  hardware.} End-to-end task latency at INT4 is roughly $2\times$ FP16
  on a T4 (6.96\,s vs.\ 3.59\,s per task on our diagnostic suite),
  because dequantization overhead dominates when compute, not memory
  bandwidth, is the bottleneck (Q2).
  \item \textbf{Fine-tune with care: short generic traces suppress
  follow-up search.} QLoRA on 15 short ReAct demonstrations leaves
  diagnostic-suite accuracy unchanged (98.0\%) but degrades the standard
  benchmarks (pooled McNemar $p{=}0.041$), with the effect concentrated
  on HotpotQA. The behavioral signal is far stronger than the accuracy
  signal: mean executed searches per question fall from 1.48 to 1.13, and
  the proportion of questions receiving a second search collapses from 39\%
  to 10\% (Wilcoxon signed-rank $p{=}8.8\times10^{-6}$, rank-biserial
  $r{=}0.83$) (Q3).
\end{enumerate}

We view the third finding---a well-diagnosed negative result---as the
most actionable: it identifies a concrete failure mode that any
practitioner fine-tuning small agent models on demonstration traces
should measure, and that per-benchmark accuracy alone does not reveal.

\section{Related Work}

\paragraph{Small language models.}
Capable sub-10B models---Phi-3~\cite{abdin2024phi3},
Gemma~\cite{gemma2024}, Mistral-7B~\cite{jiang2023mistral},
Qwen2~\cite{yang2024qwen2}---have narrowed the gap to much larger models,
and surveys~\cite{wan2024efficient,lu2024slmsurvey} catalogue the design
space. \citet{belcak2025slm} argue that SLMs will dominate agentic
deployments; \citet{srivastava2026effgen} build a framework enabling SLMs as
autonomous agents. Our work is complementary: rather than proposing a
framework, we measure how the two standard efficiency levers affect agentic
behavior under a minimal scaffold.

\paragraph{Quantization and parameter-efficient fine-tuning.}
Post-training quantization to 8-bit~\cite{dettmers2022llmint8} and
4-bit~\cite{frantar2023gptq,dettmers2023qlora} preserves perplexity and
single-turn accuracy remarkably well, and LoRA~\cite{hu2022lora} and
QLoRA~\cite{dettmers2023qlora} enable fine-tuning with a fraction of a
percent of trainable parameters; distillation offers a complementary
route~\cite{gu2024minillm}. Most quantization evaluations measure static
benchmarks; we measure closed-loop agent behavior, where errors compound
across steps and the number of tool invocations is itself an outcome.

\paragraph{Language agents and benchmarks.}
ReAct~\cite{yao2023react} interleaves reasoning traces with tool calls
and remains the canonical minimal agent scaffold;
Toolformer~\cite{schick2023toolformer} showed models can learn API use
self-supervised. Agent benchmarks such as AgentBench~\cite{liu2024agentbench},
ALFWorld~\cite{shridhar2021alfworld} and WebArena~\cite{zhou2024webarena}
target large models and heavy environments, and are poorly matched to 3.8B
models on free hardware where format compliance dominates the score. We
therefore combine a controlled diagnostic suite with
GSM8K~\cite{cobbe2021gsm8k} (calculator tool) and
HotpotQA~\cite{yang2018hotpotqa} in the distractor setting (search tool).

\section{Method}
\label{sec:method}

\subsection{Agent architecture}
We implement a minimal ReAct agent. The model receives a system prompt
specifying the available tools and the
\texttt{Thought}/\texttt{Action}/\texttt{Observation} format, and
generates greedily (temperature 0, \texttt{do\_sample=False}). A stopping
criterion halts generation whenever the model emits
\texttt{Observation:}; the named tool is then \emph{actually executed},
its real output is injected into the context, and generation resumes.
This guarantees observations are never hallucinated. The loop ends when
the model emits \texttt{Final Answer:} or a step budget is reached (6
steps; 8 for GSM8K). Two implementation details proved necessary for
small models: (i) a minimum-new-tokens guard (20 tokens) before the
stopping criterion is armed, preventing spurious first-token stops, and
(ii) trace trimming that truncates generation drift (e.g., fabricated
follow-up questions) at known markers.

If no \texttt{Final Answer:} span is produced, the answer field falls
back to the last line of the trace. We audited this fallback across all
750 agent runs: 21 answers fell back to observation text, and none of
them scored correct, so the fallback does not inflate any reported
accuracy.

\subsection{Tools}
\texttt{calculator[expr]} evaluates sanitized arithmetic expressions
(rejecting any input outside a numeric-and-operator character class).
\texttt{search[query]} retrieves from a benchmark-specific source: a
fixed 27-entry fact base for the diagnostic suite, and, for HotpotQA, the
10 paragraphs (2 gold, 8 distractors) provided with each question, ranked by
title-weighted lexical overlap, returning the single highest-scoring
paragraph truncated to 600 characters. Because only one paragraph is
returned per call, any question requiring both gold paragraphs
necessarily requires at least two search calls---a fact we rely on in
Section~\ref{sec:q3}.

\subsection{Benchmarks}
\label{sec:benchmarks}
\textbf{Custom-50 (diagnostic).} 50 tasks in five categories of 10:
\emph{single\_lookup}, \emph{arithmetic}, \emph{multi\_step} (search then
calculate), \emph{tool\_selection} (must choose the correct tool), and
\emph{sequential} (explicit multi-action instructions). Scoring is
substring match of the expected value in the final answer. This suite is
a \emph{format-compliance and tool-routing check}, not a capability
measure: all configurations saturate it (Section~\ref{sec:q1}), so it
discriminates only in the sense of confirming that the agent loop
functions identically across precisions.
\textbf{GSM8K.} 100 problems sampled (seed 42) from the test split;
calculator tool only; one worked example in the system prompt; scoring is
exact numeric match against the gold answer.
\textbf{HotpotQA (distractor).} 100 validation questions with span
answers (yes/no questions excluded to keep inclusion-match scoring
unambiguous), sampled with seed 42; scoring is normalized-answer
inclusion (lowercasing, punctuation and article removal).

\subsection{Configurations}
\label{sec:configs}
\textbf{FP16}: the unmodified model in half precision (7.64\,GB weight
memory).
\textbf{INT4}: NF4 double quantization via
bitsandbytes~\cite{dettmers2023qlora} with FP16 compute (2.21\,GB).
\textbf{INT4+QLoRA}: the INT4 model with LoRA adapters ($r{=}16$,
$\alpha{=}32$) on the attention projections (\texttt{qkv\_proj},
\texttt{o\_proj}), 9{,}437{,}184 trainable parameters (0.2464\%), trained
for 8 epochs at learning rate $2\times10^{-4}$ (cosine schedule, warmup
ratio 0.1, effective batch size 4, paged 8-bit AdamW) on \textbf{15
hand-written generic ReAct demonstrations} covering the five task
patterns. The demonstrations average 1.60 tool calls each and are fully
disjoint from all evaluation tasks in entities, numbers, and wording;
\emph{no GSM8K or HotpotQA data is used}. Training takes 1.8 minutes on a
single T4.
All runs use one T4 GPU, greedy decoding, seed 42, and identical harness
code; the only variable is the model configuration. Environment:
transformers 4.41.2, torch 2.10.0+cu128, datasets 2.19.1, peft 0.11.1,
trl 0.8.6, accelerate 0.30.1, CUDA 12.8, Python 3.12, NVIDIA Tesla T4.
The bitsandbytes version for the replication and trace-depth analysis of
Section~\ref{sec:replication} is 0.50.1; for the original evaluation round
reported in Table~\ref{tab:standard} it was
installed as \texttt{>=0.46.1} and the resolved version was not recorded.
Given our own findings in Section~\ref{sec:repro} we flag this as a gap
rather than omit it; Section~\ref{sec:replication} quantifies how much it
matters in practice.

\section{Results}

\subsection{Q1: Quantization is essentially free for agentic capability}
\label{sec:q1}

On the diagnostic suite all three configurations score 98.0\% (49/50), with
the single failure in \emph{multi\_step} and every other category at 100\%.
The paired view is stronger than the aggregate: FP16 and INT4 have
\textbf{zero discordant pairs}---they succeed and fail on exactly the same 50
tasks. This is an equivalence result on this suite, not merely an undetected
difference. Mean latency is 3.59\,s (FP16), 6.96\,s (INT4) and 7.37\,s
(INT4+QLoRA).

Table~\ref{tab:standard} shows the standard benchmarks. INT4 loses 3
points on GSM8K (71\%$\rightarrow$68\%) and 3 points on HotpotQA
(47\%$\rightarrow$44\%). Paired McNemar tests find no significant
difference in either case (GSM8K: 8 vs.\ 5 discordant, $p{=}0.58$;
HotpotQA: 10 vs.\ 7 discordant, $p{=}0.63$). We state the result as a
bound rather than a null: the paired 95\% confidence intervals are
$[-4.3, +10.8]$ on GSM8K and $[-5.1, +11.1]$ on HotpotQA, so our design
rules out an INT4 penalty larger than roughly 11 points but cannot
resolve differences below that. Within that bound, quantization costs
nothing measurable while removing 71\% of weight memory.

\begin{table}[t]
\centering
\caption{Standard benchmarks (100 tasks each, seed 42). The agent is
\emph{not} fine-tuned on either dataset. Tool calls are the mean number
of executed tool invocations per task. $p$-values are exact McNemar
against the INT4 row within the same benchmark.}
\label{tab:standard}
\small
\begin{tabular}{llcccc}
\toprule
Benchmark & Config & Accuracy (\%) & Latency (s) & Tool calls & $p$ vs.\ INT4 \\
\midrule
\multirow{3}{*}{GSM8K}
 & FP16        & \textbf{71.0} & \textbf{7.32} & 3.38 & 0.581 \\
 & INT4        & 68.0          & 12.31         & 3.22 & --- \\
 & INT4+QLoRA  & 65.0          & 14.31         & 3.25 & 0.388 \\
\midrule
\multirow{3}{*}{HotpotQA (distractor)}
 & FP16        & \textbf{47.0} & \textbf{5.46} & 1.65 & 0.629 \\
 & INT4        & 44.0          & 9.03          & 1.47 & --- \\
 & INT4+QLoRA  & 38.0          & 7.30          & \textbf{1.13} & 0.070 \\
\bottomrule
\end{tabular}
\end{table}

\subsection{Q2: On a T4, quantization buys memory, not speed}

Across every benchmark, FP16 is the fastest configuration---roughly
$2\times$ faster than INT4 per task---despite INT4's smaller weights. On a T4, NF4
weights must be dequantized to FP16 for every matrix multiplication, and
this overhead dominates at batch size 1, where the GPU is not
memory-bandwidth bound. The T4 is a Turing-generation accelerator with no
native low-precision matmul path for NF4; this result is a property of
that hardware class and should not be read as a general claim about
quantized inference.

The practical implication is a deployment decision rule: choose INT4 when
the constraint is \emph{fitting} the model (edge devices, multi-tenant
serving, freeing VRAM for context), and FP16 when the constraint is
\emph{latency} and the model already fits. Reported speedups from
quantization on server-class GPUs with fused low-bit kernels do not
automatically transfer to commodity accelerators, and end-to-end agent
latency---which compounds per-step generation over multiple tool
calls---amplifies the difference.

We report weight memory (7.64\,GB vs.\ 2.21\,GB) as measured by
\texttt{get\_memory\_footprint()}. This excludes activations and the KV
cache; peak allocated memory during evaluation is the quantity that
determines whether a model fits on a given device, and our harness now
records it alongside every run.

\subsection{Q3: Generic-trace QLoRA suppresses follow-up search}
\label{sec:q3}

QLoRA fine-tuning leaves the diagnostic suite unchanged (98.0\%). On the
standard benchmarks it costs 3 points on GSM8K and 6 points on HotpotQA.
Neither benchmark alone reaches significance (GSM8K $p{=}0.39$; HotpotQA
$p{=}0.070$), but the effect is directionally consistent across both, and
pooling the discordant pairs over the two benchmarks gives 15 vs.\ 5,
exact McNemar $p{=}0.041$. The accuracy evidence is therefore suggestive
but underpowered at $n{=}100$, and we do not rest the claim on it.

The decisive evidence is behavioral. Table~\ref{tab:depth} gives the
distribution of executed tool calls per HotpotQA question. The base INT4
agent issues a second search on 39 of 100 questions; after fine-tuning it
does so on 10. Across the 100 paired questions the fine-tuned agent used
fewer calls on 30 and more on 2 (paired Wilcoxon signed-rank
$W{=}44$, $p{=}8.8\times10^{-6}$, rank-biserial $r{=}0.83$).

\begin{table}[t]
\centering
\caption{Distribution of executed tool calls per HotpotQA question
(replication run, $n{=}100$ paired). Our search tool returns one paragraph
per call, so any question needing both gold paragraphs requires at least two
calls. Fine-tuning removes three quarters of the second searches.}
\label{tab:depth}
\small
\begin{tabular}{lccccc}
\toprule
Config & 0 calls & 1 call & 2 calls & 3+ calls & $\ge$2 calls \\
\midrule
INT4          & 1 & 60 & 35 & 4 & \textbf{39} \\
INT4+QLoRA    & 1 & 89 &  9 & 1 & \textbf{10} \\
\bottomrule
\end{tabular}
\end{table}

Second searches are not idle. In both configurations, questions that
received at least two calls scored substantially higher than those that
received one (INT4: 51.3\% vs.\ 37.7\%; INT4+QLoRA: 60.0\% vs.\ 35.6\%).
Fine-tuning therefore suppresses precisely the behavior that carries the
multi-hop questions. Consistent with this, fine-tuned HotpotQA latency
\emph{drops} (8.24\,s $\rightarrow$ 6.51\,s): the agent is faster because
it gives up earlier.

Two observations constrain the mechanism, and we state them because they
rule out the simplest explanation. First, our 15 demonstrations average
1.60 tool calls, so the fine-tuned model at 1.13 has gone \emph{below}
the training distribution rather than converging to it. Second, on GSM8K
mean tool calls did not fall at all (3.22 $\rightarrow$ 3.25). The effect
is therefore not a global shortening of trace length, which would have
affected both benchmarks. What the demonstrations have in common is that
every search they contain returns a directly usable fact; none models
recovering from an unhelpful observation. We interpret the effect as a
reduced propensity to issue a \emph{follow-up} search after an
observation that does not already contain the answer---a search-specific
behavior---rather than a learned length prior.

We term this failure mode \textbf{search-horizon truncation}. It is
invisible to in-domain evaluation (where one or two calls suffice), and at
$n{=}100$ it is not reliably detectable from accuracy alone: the same
intervention that moves accuracy by an amount indistinguishable from noise
moves trace depth by five orders of magnitude in $p$. This asymmetry is the
practical lesson. We recommend that agent fine-tuning reports include
tool-call distributions, not merely means, alongside accuracy.

\paragraph{Ablation: do longer demonstrations repair it?}
We expect that mixing longer multi-hop demonstrations into the training set
repairs the effect, since the demonstrations currently contain no example of
recovering from an unhelpful observation. We leave this to future work.

\subsection{Replication under a different quantization kernel}
\label{sec:replication}

The HotpotQA INT4 and INT4+QLoRA configurations were re-run in a fresh
session, in separate processes, under bitsandbytes 0.50.1---a different
kernel version from the original round. Accuracy reproduced to within one
task (INT4 44\%$\rightarrow$43\%; INT4+QLoRA 38\%$\rightarrow$38\%) and mean
tool calls reproduced to two decimal places (1.47$\rightarrow$1.48;
1.13$\rightarrow$1.13). The trace-depth analysis above uses this replication
run, for which per-task tool-call counts and full agent traces are released.

This is a useful counterweight to Section~\ref{sec:repro}: between two recent
bitsandbytes releases under a fixed harness the agent's behavior is stable, so
the variance we observed earlier is attributable to the harness and kernel
changing together rather than to quantization kernels being inherently
unstable. We also record peak allocated memory during evaluation: 3.77\,GB for
INT4 against 2.21\,GB of weights---the gap is activations and KV cache, and it
is the peak figure that determines whether a deployment target fits.

\subsection{Reproducibility: quantized-agent scores are environment-sensitive}
\label{sec:repro}

In an earlier evaluation round conducted with an older bitsandbytes
kernel and a slightly different harness, the identical INT4 model scored
92\% on a 25-task subset of our diagnostic suite (multi\_step: 60\%),
and QLoRA appeared to ``recover'' accuracy to 96\%. After unifying the
harness and upgrading the quantization library, base INT4 scores 98\%
and the recovery effect disappears. We report this deliberately: small
quantized models operate close to capability thresholds on multi-step
tasks, so kernel-level numerical differences can move headline agentic
scores by several points. Pinned library versions and a single shared
harness across all configurations are therefore not optional hygiene but
a validity requirement for quantized-agent research.

\section{Limitations}
\label{sec:limitations}

Our study covers one model family (Phi-3-mini), one quantization method
(bitsandbytes NF4) and one hardware target (T4); we test no INT8 intermediate
point and no alternative scheme such as GPTQ~\cite{frantar2023gptq} or AWQ,
and the latency conclusions in particular will differ on GPUs with fused
low-bit kernels. The diagnostic suite is small (50 tasks), uses a closed
27-entry fact base and is saturated by every configuration, so it functions as
a format-compliance check rather than a capability measurement; differences of
$\le$3 points on the 100-task benchmarks are within our resolution, and
HotpotQA inclusion match on span answers is more permissive than exact match.

The FP16 configuration was evaluated in a separate cloud session from the INT4
configurations. Accuracy comparisons are unaffected---decoding is greedy and
the harness byte-identical---but the latency comparison is not concurrent and
shared-tenancy GPUs vary with host load.

The QLoRA study uses a single minimal demonstration set (15 traces) and a
single training seed; we show that this regime suppresses follow-up search,
not that QLoRA cannot help agents trained on longer or task-matched traces.
Our INT4 and INT4+QLoRA configurations also differ by the dtype changes that
\texttt{prepare\_model\_for\_kbit\_training} introduces (layer norms and the LM
head are upcast to fp32) as well as by the adapters. We report no zero-adapter
control, so attribution to the trained weights rests on the tool-call
mechanism rather than an isolated ablation; our released harness implements
this control. Finally, the scaffold is intentionally minimal (ReAct, greedy
decoding, two tools) and results may differ under richer scaffolds with
reflection or planning.

\section{Broader impacts}
\label{sec:impacts}

Reducing the memory cost of capable tool-using agents supports on-device and
on-premise deployment, where data need not leave the user's hardware, and
lowers the cost of agentic inference; the same accessibility is dual-use,
since cheaper agents are also cheaper to run at scale for unwanted automation.
Our negative result is itself safety-relevant: we document a cheap fine-tuning
regime that leaves headline accuracy close to unchanged while measurably
suppressing an agent's willingness to gather evidence before answering. An
agent that stops searching early is not obviously broken from the outside; it
answers more confidently on less information. Practitioners adapting small
agents on demonstration data should measure tool-call depth, not accuracy
alone.

\section{Conclusion}

A 3.8B-parameter model quantized to 4 bits retains essentially all of its
tool-using agentic capability---identical per-task outcomes on our diagnostic
suite, no statistically detectable difference on two standard
benchmarks---at 29\% of the weight memory. Two caveats qualify the picture:
on commodity GPUs quantization trades latency for memory rather than improving
both, and lightweight fine-tuning on short demonstration traces can silently
suppress follow-up tool use, degrading exactly the multi-hop behavior agents
exist to perform. Measuring trace depth, pinning the numerical environment and using paired
tests should be standard practice for efficient-agent research.

\paragraph{Reproducibility.} Code, harness, demonstrations, per-task CSVs and
the analysis script are anonymized at \url{ANONYMIZED-REPO-URL}.
% TODO(anonymous): paste an anonymous.4open.science link. Do NOT paste a
% GitHub URL containing your username, and check the mirrored notebooks for
% identifying strings (account names, home paths) before publishing it.



\bibliographystyle{plainnat}
\begin{thebibliography}{Srivastava et~al.(2026)}

\bibitem[Belcak and Heinrich(2025)]{belcak2025slm}
P.~Belcak and G.~Heinrich.
\newblock Small language models are the future of agentic {AI}.
\newblock \emph{arXiv preprint arXiv:2506.02153}, 2025.

\bibitem[Srivastava et~al.(2026)]{srivastava2026effgen}
G.~Srivastava, A.~Hussain, C.~Wang, Y.~C. Lin, and X.~Wang.
\newblock {EffGen}: Enabling small language models as capable autonomous
  agents.
\newblock In \emph{ICML}, 2026. arXiv:2602.00887.

\bibitem{abdin2024phi3}
M.~Abdin et~al.
\newblock Phi-3 technical report: A highly capable language model locally
  on your phone.
\newblock \emph{arXiv preprint arXiv:2404.14219}, 2024.

\bibitem{yao2023react}
S.~Yao, J.~Zhao, D.~Yu, N.~Du, I.~Shafran, K.~Narasimhan, and Y.~Cao.
\newblock {ReAct}: Synergizing reasoning and acting in language models.
\newblock In \emph{ICLR}, 2023.

\bibitem{dettmers2023qlora}
T.~Dettmers, A.~Pagnoni, A.~Holtzman, and L.~Zettlemoyer.
\newblock {QLoRA}: Efficient finetuning of quantized {LLM}s.
\newblock In \emph{NeurIPS}, 2023.

\bibitem{hu2022lora}
E.~J. Hu, Y.~Shen, P.~Wallis, Z.~Allen-Zhu, Y.~Li, S.~Wang, L.~Wang, and
  W.~Chen.
\newblock {LoRA}: Low-rank adaptation of large language models.
\newblock In \emph{ICLR}, 2022.

\bibitem{dettmers2022llmint8}
T.~Dettmers, M.~Lewis, Y.~Belkada, and L.~Zettlemoyer.
\newblock {LLM.int8()}: 8-bit matrix multiplication for transformers at
  scale.
\newblock In \emph{NeurIPS}, 2022.

\bibitem{frantar2023gptq}
E.~Frantar, S.~Ashkboos, T.~Hoefler, and D.~Alistarh.
\newblock {GPTQ}: Accurate post-training quantization for generative
  pre-trained transformers.
\newblock In \emph{ICLR}, 2023.

\bibitem{gu2024minillm}
Y.~Gu, L.~Dong, F.~Wei, and M.~Huang.
\newblock {MiniLLM}: Knowledge distillation of large language models.
\newblock In \emph{ICLR}, 2024.

\bibitem{schick2023toolformer}
T.~Schick, J.~Dwivedi-Yu, R.~Dess{\`i}, R.~Raileanu, M.~Lomeli,
  L.~Zettlemoyer, N.~Cancedda, and T.~Scialom.
\newblock Toolformer: Language models can teach themselves to use tools.
\newblock In \emph{NeurIPS}, 2023.

\bibitem{liu2024agentbench}
X.~Liu et~al.
\newblock {AgentBench}: Evaluating {LLM}s as agents.
\newblock In \emph{ICLR}, 2024.

\bibitem{shridhar2021alfworld}
M.~Shridhar, X.~Yuan, M.-A. C{\^o}t{\'e}, Y.~Bisk, A.~Trischler, and
  M.~Hausknecht.
\newblock {ALFWorld}: Aligning text and embodied environments for
  interactive learning.
\newblock In \emph{ICLR}, 2021.

\bibitem{zhou2024webarena}
S.~Zhou et~al.
\newblock {WebArena}: A realistic web environment for building autonomous
  agents.
\newblock In \emph{ICLR}, 2024.

\bibitem{cobbe2021gsm8k}
K.~Cobbe et~al.
\newblock Training verifiers to solve math word problems.
\newblock \emph{arXiv preprint arXiv:2110.14168}, 2021.

\bibitem{yang2018hotpotqa}
Z.~Yang, P.~Qi, S.~Zhang, Y.~Bengio, W.~W. Cohen, R.~Salakhutdinov, and
  C.~D. Manning.
\newblock {HotpotQA}: A dataset for diverse, explainable multi-hop
  question answering.
\newblock In \emph{EMNLP}, 2018.

\bibitem{gemma2024}
Gemma Team.
\newblock Gemma: Open models based on {Gemini} research and technology.
\newblock \emph{arXiv preprint arXiv:2403.08295}, 2024.

\bibitem{jiang2023mistral}
A.~Q. Jiang et~al.
\newblock Mistral 7{B}.
\newblock \emph{arXiv preprint arXiv:2310.06825}, 2023.

\bibitem{yang2024qwen2}
A.~Yang et~al.
\newblock Qwen2 technical report.
\newblock \emph{arXiv preprint arXiv:2407.10671}, 2024.

\bibitem{wan2024efficient}
Z.~Wan et~al.
\newblock Efficient large language models: A survey.
\newblock \emph{Transactions on Machine Learning Research}, 2024.

\bibitem{lu2024slmsurvey}
Z.~Lu et~al.
\newblock Small language models: Survey, measurements, and insights.
\newblock \emph{arXiv preprint arXiv:2409.15790}, 2024.

\end{thebibliography}

\newpage
\section*{NeurIPS Paper Checklist}




\begin{enumerate}

\item {\bf Claims}
    \item[] Question: Do the main claims made in the abstract and introduction accurately reflect the paper's contributions and scope?
    \item[] Answer: \answerYes{} % Replace by \answerYes{}, \answerNo{}, or \answerNA{}.
    \item[] Justification: The abstract and Section~\ref{sec:intro} state three claims that map one-to-one
onto Sections~\ref{sec:q1}--\ref{sec:q3}. Scope is stated explicitly: one model
(Phi-3-mini), one quantization method (bitsandbytes NF4), one hardware target
(T4), and one agent scaffold (ReAct). The quantization result is stated as a
bound rather than an equivalence where the sample size supports only a bound,
and we state explicitly that the accuracy evidence for the fine-tuning result
is underpowered and that the claim rests on the behavioral measure instead.
    \item[] Guidelines:
    \begin{itemize}
        \item The answer \answerNA{} means that the abstract and introduction do not include the claims made in the paper.
        \item The abstract and/or introduction should clearly state the claims made, including the contributions made in the paper and important assumptions and limitations. A \answerNo{} or \answerNA{} answer to this question will not be perceived well by the reviewers. 
        \item The claims made should match theoretical and experimental results, and reflect how much the results can be expected to generalize to other settings. 
        \item It is fine to include aspirational goals as motivation as long as it is clear that these goals are not attained by the paper. 
    \end{itemize}

\item {\bf Limitations}
    \item[] Question: Does the paper discuss the limitations of the work performed by the authors?
    \item[] Answer: \answerYes{} % Replace by \answerYes{}, \answerNo{}, or \answerNA{}.
    \item[] Justification: Section~\ref{sec:limitations} is devoted to limitations: the single model
family, the absence of an INT8 intermediate condition and of alternative
quantization schemes, the saturation of the diagnostic suite, the resolution
floor at $n{=}100$, the permissiveness of inclusion-match scoring, the fact
that the FP16 configuration was evaluated in a separate cloud session so
latency comparisons are not concurrent, the single training seed, the absence
of a zero-adapter control separating adapter effects from preparation-induced
dtype changes, and the minimality of the scaffold.
    \item[] Guidelines:
    \begin{itemize}
        \item The answer \answerNA{} means that the paper has no limitation while the answer \answerNo{} means that the paper has limitations, but those are not discussed in the paper. 
        \item The authors are encouraged to create a separate ``Limitations'' section in their paper.
        \item The paper should point out any strong assumptions and how robust the results are to violations of these assumptions (e.g., independence assumptions, noiseless settings, model well-specification, asymptotic approximations only holding locally). The authors should reflect on how these assumptions might be violated in practice and what the implications would be.
        \item The authors should reflect on the scope of the claims made, e.g., if the approach was only tested on a few datasets or with a few runs. In general, empirical results often depend on implicit assumptions, which should be articulated.
        \item The authors should reflect on the factors that influence the performance of the approach. For example, a facial recognition algorithm may perform poorly when image resolution is low or images are taken in low lighting. Or a speech-to-text system might not be used reliably to provide closed captions for online lectures because it fails to handle technical jargon.
        \item The authors should discuss the computational efficiency of the proposed algorithms and how they scale with dataset size.
        \item If applicable, the authors should discuss possible limitations of their approach to address problems of privacy and fairness.
        \item While the authors might fear that complete honesty about limitations might be used by reviewers as grounds for rejection, a worse outcome might be that reviewers discover limitations that aren't acknowledged in the paper. The authors should use their best judgment and recognize that individual actions in favor of transparency play an important role in developing norms that preserve the integrity of the community. Reviewers will be specifically instructed to not penalize honesty concerning limitations.
    \end{itemize}

\item {\bf Theory assumptions and proofs}
    \item[] Question: For each theoretical result, does the paper provide the full set of assumptions and a complete (and correct) proof?
    \item[] Answer: \answerNA{} % Replace by \answerYes{}, \answerNo{}, or \answerNA{}.
    \item[] Justification: The paper is purely empirical and contains no theoretical results.
    \item[] Guidelines:
    \begin{itemize}
        \item The answer \answerNA{} means that the paper does not include theoretical results. 
        \item All the theorems, formulas, and proofs in the paper should be numbered and cross-referenced.
        \item All assumptions should be clearly stated or referenced in the statement of any theorems.
        \item The proofs can either appear in the main paper or the supplemental material, but if they appear in the supplemental material, the authors are encouraged to provide a short proof sketch to provide intuition. 
        \item Inversely, any informal proof provided in the core of the paper should be complemented by formal proofs provided in appendix or supplemental material.
        \item Theorems and Lemmas that the proof relies upon should be properly referenced. 
    \end{itemize}

    \item {\bf Experimental result reproducibility}
    \item[] Question: Does the paper fully disclose all the information needed to reproduce the main experimental results of the paper to the extent that it affects the main claims and/or conclusions of the paper (regardless of whether the code and data are provided or not)?
    \item[] Answer: \answerYes{} % Replace by \answerYes{}, \answerNo{}, or \answerNA{}.
    \item[] Justification: Section~\ref{sec:method} specifies the agent loop, stopping criterion, tool
implementations, step budgets, decoding settings, benchmark sampling seeds,
scoring functions, quantization configuration, LoRA hyperparameters, and the
full training configuration. Library versions are pinned in
Section~\ref{sec:configs}, including the bitsandbytes version for the
replication run; where a version was not recorded for the earlier round we say
so rather than omit it. The released repository contains the harness, the 15
demonstrations, per-task results, and a script that regenerates every reported
statistic.
    \item[] Guidelines:
    \begin{itemize}
        \item The answer \answerNA{} means that the paper does not include experiments.
        \item If the paper includes experiments, a \answerNo{} answer to this question will not be perceived well by the reviewers: Making the paper reproducible is important, regardless of whether the code and data are provided or not.
        \item If the contribution is a dataset and\slash or model, the authors should describe the steps taken to make their results reproducible or verifiable. 
        \item Depending on the contribution, reproducibility can be accomplished in various ways. For example, if the contribution is a novel architecture, describing the architecture fully might suffice, or if the contribution is a specific model and empirical evaluation, it may be necessary to either make it possible for others to replicate the model with the same dataset, or provide access to the model. In general. releasing code and data is often one good way to accomplish this, but reproducibility can also be provided via detailed instructions for how to replicate the results, access to a hosted model (e.g., in the case of a large language model), releasing of a model checkpoint, or other means that are appropriate to the research performed.
        \item While NeurIPS does not require releasing code, the conference does require all submissions to provide some reasonable avenue for reproducibility, which may depend on the nature of the contribution. For example
        \begin{enumerate}
            \item If the contribution is primarily a new algorithm, the paper should make it clear how to reproduce that algorithm.
            \item If the contribution is primarily a new model architecture, the paper should describe the architecture clearly and fully.
            \item If the contribution is a new model (e.g., a large language model), then there should either be a way to access this model for reproducing the results or a way to reproduce the model (e.g., with an open-source dataset or instructions for how to construct the dataset).
            \item We recognize that reproducibility may be tricky in some cases, in which case authors are welcome to describe the particular way they provide for reproducibility. In the case of closed-source models, it may be that access to the model is limited in some way (e.g., to registered users), but it should be possible for other researchers to have some path to reproducing or verifying the results.
        \end{enumerate}
    \end{itemize}


\item {\bf Open access to data and code}
    \item[] Question: Does the paper provide open access to the data and code, with sufficient instructions to faithfully reproduce the main experimental results, as described in supplemental material?
    \item[] Answer: \answerYes{} % Replace by \answerYes{}, \answerNo{}, or \answerNA{}.
    \item[] Justification: An anonymized repository is provided for review. It contains the evaluation
harness, model-loading code for every configuration, the training script, the
15 demonstrations, per-task CSVs including tool-call counts and complete agent
traces, the trained adapters, captured environment metadata, and a single
script reproducing the full pipeline. Benchmarks derive from public sources
(GSM8K, HotpotQA) or are released in full (the diagnostic suite and its fact
base). One results directory was recovered from notebook output after a cloud
session was lost; it is marked as such, its provenance is documented, and the
source notebooks are included with outputs intact so the recovery is
auditable.
    \item[] Guidelines:
    \begin{itemize}
        \item The answer \answerNA{} means that paper does not include experiments requiring code.
        \item Please see the NeurIPS code and data submission guidelines (\url{https://neurips.cc/public/guides/CodeSubmissionPolicy}) for more details.
        \item While we encourage the release of code and data, we understand that this might not be possible, so \answerNo{} is an acceptable answer. Papers cannot be rejected simply for not including code, unless this is central to the contribution (e.g., for a new open-source benchmark).
        \item The instructions should contain the exact command and environment needed to run to reproduce the results. See the NeurIPS code and data submission guidelines (\url{https://neurips.cc/public/guides/CodeSubmissionPolicy}) for more details.
        \item The authors should provide instructions on data access and preparation, including how to access the raw data, preprocessed data, intermediate data, and generated data, etc.
        \item The authors should provide scripts to reproduce all experimental results for the new proposed method and baselines. If only a subset of experiments are reproducible, they should state which ones are omitted from the script and why.
        \item At submission time, to preserve anonymity, the authors should release anonymized versions (if applicable).
        \item Providing as much information as possible in supplemental material (appended to the paper) is recommended, but including URLs to data and code is permitted.
    \end{itemize}


\item {\bf Experimental setting/details}
    \item[] Question: Does the paper specify all the training and test details (e.g., data splits, hyperparameters, how they were chosen, type of optimizer) necessary to understand the results?
    \item[] Answer: \answerYes{} % Replace by \answerYes{}, \answerNo{}, or \answerNA{}.
    \item[] Justification: Section~\ref{sec:configs} gives the LoRA rank and scaling, target modules,
trainable-parameter count, epochs, learning rate, schedule, warmup ratio,
effective batch size, optimizer, sequence length, seed, and wall-clock training
time. Section~\ref{sec:benchmarks} gives benchmark sizes, sampling seeds,
exclusion criteria, and scoring rules. Decoding is greedy throughout.
    \item[] Guidelines:
    \begin{itemize}
        \item The answer \answerNA{} means that the paper does not include experiments.
        \item The experimental setting should be presented in the core of the paper to a level of detail that is necessary to appreciate the results and make sense of them.
        \item The full details can be provided either with the code, in appendix, or as supplemental material.
    \end{itemize}

\item {\bf Experiment statistical significance}
    \item[] Question: Does the paper report error bars suitably and correctly defined or other appropriate information about the statistical significance of the experiments?
    \item[] Answer: \answerYes{} % Replace by \answerYes{}, \answerNo{}, or \answerNA{}.
    \item[] Justification: Because every configuration runs the same task list in the same order, results
are paired, and we use McNemar's exact test for accuracy and Wilcoxon
signed-rank for tool-call depth rather than comparing independent proportions.
We report discordant-pair counts, paired difference confidence intervals, exact
$p$-values, and a rank-biserial effect size. We state where results are not
significant, including for our own fine-tuning result on accuracy, and we do
not claim equivalence where the interval supports only a bound. The
single-seed limitation is disclosed in Section~\ref{sec:limitations}.
    \item[] Guidelines:
    \begin{itemize}
        \item The answer \answerNA{} means that the paper does not include experiments.
        \item The authors should answer \answerYes{} if the results are accompanied by error bars, confidence intervals, or statistical significance tests, at least for the experiments that support the main claims of the paper.
        \item The factors of variability that the error bars are capturing should be clearly stated (for example, train/test split, initialization, random drawing of some parameter, or overall run with given experimental conditions).
        \item The method for calculating the error bars should be explained (closed form formula, call to a library function, bootstrap, etc.)
        \item The assumptions made should be given (e.g., Normally distributed errors).
        \item It should be clear whether the error bar is the standard deviation or the standard error of the mean.
        \item It is OK to report 1-sigma error bars, but one should state it. The authors should preferably report a 2-sigma error bar than state that they have a 96\% CI, if the hypothesis of Normality of errors is not verified.
        \item For asymmetric distributions, the authors should be careful not to show in tables or figures symmetric error bars that would yield results that are out of range (e.g., negative error rates).
        \item If error bars are reported in tables or plots, the authors should explain in the text how they were calculated and reference the corresponding figures or tables in the text.
    \end{itemize}

\item {\bf Experiments compute resources}
    \item[] Question: For each experiment, does the paper provide sufficient information on the computer resources (type of compute workers, memory, time of execution) needed to reproduce the experiments?
    \item[] Answer: \answerYes{} % Replace by \answerYes{}, \answerNo{}, or \answerNA{}.
    \item[] Justification: All experiments run on a single NVIDIA T4 on a free cloud tier. We report
weight memory (7.64\,GB FP16, 2.21\,GB INT4), peak allocated memory during
evaluation (3.77\,GB INT4), per-task latency for every configuration and
benchmark, and training wall-clock time (1.8--2.0 minutes). Total compute is a
few GPU-hours, including runs not reported in the paper.
    \item[] Guidelines:
    \begin{itemize}
        \item The answer \answerNA{} means that the paper does not include experiments.
        \item The paper should indicate the type of compute workers CPU or GPU, internal cluster, or cloud provider, including relevant memory and storage.
        \item The paper should provide the amount of compute required for each of the individual experimental runs as well as estimate the total compute. 
        \item The paper should disclose whether the full research project required more compute than the experiments reported in the paper (e.g., preliminary or failed experiments that didn't make it into the paper). 
    \end{itemize}
    
\item {\bf Code of ethics}
    \item[] Question: Does the research conducted in the paper conform, in every respect, with the NeurIPS Code of Ethics \url{https://neurips.cc/public/EthicsGuidelines}?
    \item[] Answer: \answerYes{} % Replace by \answerYes{}, \answerNo{}, or \answerNA{}.
    \item[] Justification: The work uses publicly released models and public benchmarks, involves no
human subjects and no personally identifying data, and is conducted and
reported in conformance with the NeurIPS Code of Ethics. Anonymity is preserved
in this submission.
    \item[] Guidelines:
    \begin{itemize}
        \item The answer \answerNA{} means that the authors have not reviewed the NeurIPS Code of Ethics.
        \item If the authors answer \answerNo, they should explain the special circumstances that require a deviation from the Code of Ethics.
        \item The authors should make sure to preserve anonymity (e.g., if there is a special consideration due to laws or regulations in their jurisdiction).
    \end{itemize}


\item {\bf Broader impacts}
    \item[] Question: Does the paper discuss both potential positive societal impacts and negative societal impacts of the work performed?
    \item[] Answer: \answerYes{} % Replace by \answerYes{}, \answerNo{}, or \answerNA{}.
    \item[] Justification: Discussed in Section~\ref{sec:impacts}. Reducing the hardware and energy cost
of capable tool-using agents supports on-device and on-premise deployment, with
associated privacy and accessibility benefits and reduced reliance on large
centrally-hosted models; the corresponding risk is that cheaper autonomous
agents are also cheaper to deploy at scale for harmful purposes. Our central
negative result is itself safety-relevant: it documents a fine-tuning regime
that silently degrades an agent's exploratory behavior while leaving headline
accuracy apparently intact.
    \item[] Guidelines:
    \begin{itemize}
        \item The answer \answerNA{} means that there is no societal impact of the work performed.
        \item If the authors answer \answerNA{} or \answerNo, they should explain why their work has no societal impact or why the paper does not address societal impact.
        \item Examples of negative societal impacts include potential malicious or unintended uses (e.g., disinformation, generating fake profiles, surveillance), fairness considerations (e.g., deployment of technologies that could make decisions that unfairly impact specific groups), privacy considerations, and security considerations.
        \item The conference expects that many papers will be foundational research and not tied to particular applications, let alone deployments. However, if there is a direct path to any negative applications, the authors should point it out. For example, it is legitimate to point out that an improvement in the quality of generative models could be used to generate Deepfakes for disinformation. On the other hand, it is not needed to point out that a generic algorithm for optimizing neural networks could enable people to train models that generate Deepfakes faster.
        \item The authors should consider possible harms that could arise when the technology is being used as intended and functioning correctly, harms that could arise when the technology is being used as intended but gives incorrect results, and harms following from (intentional or unintentional) misuse of the technology.
        \item If there are negative societal impacts, the authors could also discuss possible mitigation strategies (e.g., gated release of models, providing defenses in addition to attacks, mechanisms for monitoring misuse, mechanisms to monitor how a system learns from feedback over time, improving the efficiency and accessibility of ML).
    \end{itemize}
    
\item {\bf Safeguards}
    \item[] Question: Does the paper describe safeguards that have been put in place for responsible release of data or models that have a high risk for misuse (e.g., pre-trained language models, image generators, or scraped datasets)?
    \item[] Answer: \answerNA{} % Replace by \answerYes{}, \answerNo{}, or \answerNA{}.
    \item[] Justification: We release an evaluation harness, a small set of generic tool-use
demonstrations, and LoRA adapters for a publicly available model. None of these
carries a high risk of misuse.
    \item[] Guidelines:
    \begin{itemize}
        \item The answer \answerNA{} means that the paper poses no such risks.
        \item Released models that have a high risk for misuse or dual-use should be released with necessary safeguards to allow for controlled use of the model, for example by requiring that users adhere to usage guidelines or restrictions to access the model or implementing safety filters. 
        \item Datasets that have been scraped from the Internet could pose safety risks. The authors should describe how they avoided releasing unsafe images.
        \item We recognize that providing effective safeguards is challenging, and many papers do not require this, but we encourage authors to take this into account and make a best faith effort.
    \end{itemize}

\item {\bf Licenses for existing assets}
    \item[] Question: Are the creators or original owners of assets (e.g., code, data, models), used in the paper, properly credited and are the license and terms of use explicitly mentioned and properly respected?
    \item[] Answer: \answerYes{} % Replace by \answerYes{}, \answerNo{}, or \answerNA{}.
    \item[] Justification: Phi-3-mini, GSM8K, HotpotQA, ReAct, LoRA, QLoRA and bitsandbytes are cited in
the text. All assets are used under their public licenses and for their
intended research purpose.
    \item[] Guidelines:
    \begin{itemize}
        \item The answer \answerNA{} means that the paper does not use existing assets.
        \item The authors should cite the original paper that produced the code package or dataset.
        \item The authors should state which version of the asset is used and, if possible, include a URL.
        \item The name of the license (e.g., CC-BY 4.0) should be included for each asset.
        \item For scraped data from a particular source (e.g., website), the copyright and terms of service of that source should be provided.
        \item If assets are released, the license, copyright information, and terms of use in the package should be provided. For popular datasets, \url{paperswithcode.com/datasets} has curated licenses for some datasets. Their licensing guide can help determine the license of a dataset.
        \item For existing datasets that are re-packaged, both the original license and the license of the derived asset (if it has changed) should be provided.
        \item If this information is not available online, the authors are encouraged to reach out to the asset's creators.
    \end{itemize}

\item {\bf New assets}
    \item[] Question: Are new assets introduced in the paper well documented and is the documentation provided alongside the assets?
    \item[] Answer: \answerYes{} % Replace by \answerYes{}, \answerNo{}, or \answerNA{}.
    \item[] Justification: We release a 50-task diagnostic agentic benchmark with its fact base, 15 ReAct
demonstrations, an evaluation harness, trained LoRA adapters, and per-task
results. The repository documents installation, per-configuration commands,
full reproduction, file layout, methodological notes, and the provenance and
known gaps of each results directory.
    \item[] Guidelines:
    \begin{itemize}
        \item The answer \answerNA{} means that the paper does not release new assets.
        \item Researchers should communicate the details of the dataset\slash code\slash model as part of their submissions via structured templates. This includes details about training, license, limitations, etc. 
        \item The paper should discuss whether and how consent was obtained from people whose asset is used.
        \item At submission time, remember to anonymize your assets (if applicable). You can either create an anonymized URL or include an anonymized zip file.
    \end{itemize}

\item {\bf Crowdsourcing and research with human subjects}
    \item[] Question: For crowdsourcing experiments and research with human subjects, does the paper include the full text of instructions given to participants and screenshots, if applicable, as well as details about compensation (if any)? 
    \item[] Answer: \answerNA{} % Replace by \answerYes{}, \answerNo{}, or \answerNA{}.
    \item[] Justification: The paper involves no crowdsourcing and no research with human subjects.
    \item[] Guidelines:
    \begin{itemize}
        \item The answer \answerNA{} means that the paper does not involve crowdsourcing nor research with human subjects.
        \item Including this information in the supplemental material is fine, but if the main contribution of the paper involves human subjects, then as much detail as possible should be included in the main paper. 
        \item According to the NeurIPS Code of Ethics, workers involved in data collection, curation, or other labor should be paid at least the minimum wage in the country of the data collector. 
    \end{itemize}

\item {\bf Institutional review board (IRB) approvals or equivalent for research with human subjects}
    \item[] Question: Does the paper describe potential risks incurred by study participants, whether such risks were disclosed to the subjects, and whether Institutional Review Board (IRB) approvals (or an equivalent approval/review based on the requirements of your country or institution) were obtained?
    \item[] Answer: \answerNA{} % Replace by \answerYes{}, \answerNo{}, or \answerNA{}.
    \item[] Justification: The paper involves no research with human subjects.
    \item[] Guidelines:
    \begin{itemize}
        \item The answer \answerNA{} means that the paper does not involve crowdsourcing nor research with human subjects.
        \item Depending on the country in which research is conducted, IRB approval (or equivalent) may be required for any human subjects research. If you obtained IRB approval, you should clearly state this in the paper. 
        \item We recognize that the procedures for this may vary significantly between institutions and locations, and we expect authors to adhere to the NeurIPS Code of Ethics and the guidelines for their institution. 
        \item For initial submissions, do not include any information that would break anonymity (if applicable), such as the institution conducting the review.
    \end{itemize}

\item {\bf Declaration of LLM usage}
    \item[] Question: Does the paper describe the usage of LLMs if it is an important, original, or non-standard component of the core methods in this research? Note that if the LLM is used only for writing, editing, or formatting purposes and does \emph{not} impact the core methodology, scientific rigor, or originality of the research, declaration is not required.
    %this research? 
    \item[] Answer: \answerYes{} % Replace by \answerYes{}, \answerNo{}, or \answerNA{}.
    \item[] Justification: Language models are the object of study rather than a writing aid: Phi-3-mini
is the system under evaluation throughout, and its size, quantization, LoRA
fine-tuning and agent scaffold are described in full in
Section~\ref{sec:method}.
    \item[] Guidelines:
    \begin{itemize}
        \item The answer \answerNA{} means that the core method development in this research does not involve LLMs as any important, original, or non-standard components.
        \item Please refer to our LLM policy in the NeurIPS handbook for what should or should not be described.
    \end{itemize}

\end{enumerate}

\end{document}
